\section{Experiments}
\label{sec:experiments}
% PLAN (about 2 pages and one table). Design only: no result appears here. Everything below restates SAP.md / SAP_v2.md, which were frozen before the
% locked evaluations (tools/freeze.py, tools/freeze_v2.py); if this text and SAP_v2.md differ, SAP_v2.md wins and this text is corrected.
%   6.1 Models and comparators (ours, no-twin control, v1, Qure released weights, identity, CLAHE, homomorphic filtering).
%   6.2 Arms (Table arms) and the dose-response over the twin weight.
%   6.3 Protocols per hypothesis H1'-H5' (what is computed, on which set, with which test).
%   6.4 Metrics. 6.5 Downstream protocol. 6.6 Encoder screen (selection rule fixed before the screen).
\subsection{Models and comparators}
\emph{Ours} is the ensemble of three ConvNeXt-V2-T UNet++/CBAM suppressors (seeds 0--2) trained on the Qure projections with twin weight $w=16$
(Section~\ref{sec:method_v2}). It replaced the SwinV2-T configuration of the analysis plan under the pre-declared replacement rule (R2) before the
freeze. The screening rule R1 (Section~\ref{sec:encoders}) only decided which encoders were trained at the full budget; RAD-DINO, which had
missed R1 by 0.05\,dB, was trained at the full budget as an exploratory encoder, and R2 applies to every full-budget candidate. At two seeds
(development data) RAD-DINO reached $A=0.046$ and bone-region PSNR$_{cal}=21.92$\,dB and ConvNeXt-V2-T $A=0.050$ and $22.01$\,dB; both met
R2, and the pre-declared tie-break (mean absorption within 0.005, then the higher bone-region PSNR) chose ConvNeXt-V2-T, the only encoder that
passed both R1 and R2, so excluding RAD-DINO would not have changed the choice. With the third seeds the absorption gap is 0.0055, just outside
the band; the choice was not reopened, and both encoders are reported. The SwinV2-T and RAD-DINO pairs (twin and no twin, three seeds each) are secondary.
Comparators are: the \emph{no-twin control} (same encoder, data and recipe, $w=0$, three seeds); the v1 model trained on LIDC-IDRI only (\texttt{clia},
pre-specified pick of the v1 study); the publicly released weights of \citet{angaitkar2026anatomy} (SwinV2-B at $1024$~px; research licence);
and three controls that need no training: identity, CLAHE \citep{zuiderveld1994clahe} and a homomorphic filter
\citep{oppenheim1968homomorphic}. All methods are scored with the same code on the same images.

\subsection{Arms}
Table~\ref{tab:arms} lists the arms. Each arm changes one factor of its parent. The confirmatory dose-response varies the twin weight over
$w\in\{0,4,8,16,32,64\}$ in the SwinV2-T family with 3, 2, 2, 4, 2 and 2 seeds ($w=1$ appears only in the screening). The insertion-only arm and the
last four rows form the ablation of the regulariser (two seeds each, SwinV2-T, $w=16$).

\begin{table}[htbp]
\centering\small
\renewcommand{\arraystretch}{1.15}
\caption{Training arms. ``Lesions'' is the maximum number of lesions inserted into the base image (counterfactual insertion; before the physics
for the LIDC renders, before the tone curve for the Qure projections); the twin arms paste 1--2 further lesions per twin (up to 4 in
\texttt{qp\_twin16\_n4}) into the finished composite (Section~\ref{sec:method_v2});
``Twin'' is the twin weight $w$, varied over $\{4,8,16,32,64\}$ in the \texttt{qp\_twin}$w$ arms ($w=1$ in the screening only). The parent of an arm is \texttt{recipe} for the
LIDC arms, \texttt{qp} for the first Qure arms and \texttt{qp\_twin16} for the last four rows.}
\label{tab:arms}
\begin{tabular}{@{}p{0.27\linewidth}p{0.09\linewidth}p{0.10\linewidth}p{0.09\linewidth}p{0.28\linewidth}@{}}
\toprule
Arm & Source & Lesions & Twin & Change \\
\midrule
\texttt{recipe} & LIDC & 0 & none & v1 baseline \\
\texttt{nodule\_loss} & LIDC & 0 & none & lesion-weighted loss \\
\texttt{clia} & LIDC & 3 & none & counterfactual insertion \\
\texttt{clia\_nodule} & LIDC & 3 & none & insertion and lesion-weighted loss \\
\texttt{qp} & Qure & 0 & none & Qure projections (no-twin control) \\
\texttt{qp\_clia} & Qure & 3 & none & insertion only, no tie term \\
\texttt{qp\_twin}$w$ & Qure & 3 & varied & twin regulariser \\
\texttt{qp\_twin16\_sg} & Qure & 3 & 16 & stop-gradient on the un-pasted branch \\
\texttt{qp\_twin16\_wide} & Qure & 3 & 16 & support dilated with a $31\times31$ window \\
\texttt{qp\_twin16\_all} & Qure & 3 & 16 & whole-image support \\
\texttt{qp\_twin16\_n4} & Qure & 3 & 16 & four lesions per twin \\
\bottomrule
\end{tabular}
\end{table}

\subsection{Protocols per hypothesis}
The analysis plan (SAP.md for the v1 family, SAP\_v2.md for v2) was frozen before any held-out or locked evaluation, on 1~November 2026
(\path{freeze_v2.json}: hashes of both plans, the SZCH development and held-out lists, the split files and the 52 weight files, at commit
\texttt{dd52b6a}). Paired differences are tested with the Wilcoxon signed-rank test \citep{wilcoxon1945individual} and given percentile bootstrap intervals
\citep{efron1993bootstrap}. Multiplicity is handled by the Holm procedure \citep{holm1979simple} on the $p$-values within each family: with
the $m$ $p$-values of a family ordered as $p_{(1)}\le\dots\le p_{(m)}$,
\begin{equation}
\tilde p_{(j)}=\max_{i\le j}\,\min\big\{1,\,(m-i+1)\,p_{(i)}\big\}.
\label{eq:holm}
\end{equation}
Confidence intervals are unadjusted 95\,\% intervals; effect sizes are always reported with them, and a hypothesis that fails is reported as failed. The external validation on
BDCXR-3257 has its own addendum, frozen before any model saw those images (Section~\ref{sec:bdcxr_protocol}).
\begin{description}
\item[H1$'$, dose-response.] On VinDr-RibCXR (245 images, six sites each, lesions from LIDC test CTs, $\alpha=1$), the per-site absorption $A$ averaged over
 seeds falls with the twin weight: the slope of $A$ on $\log_2(1+w)$, with an image-cluster bootstrap interval, passes if its upper bound is below zero.
 The same analysis on the Qure-projection test split (exact truth) is secondary.
\item[H2$'$, fidelity.] On the 51 held-out SZCH pairs, bone-region PSNR of ours against the reference model is non-inferior within a margin of 0.5~dB:
 with the paired mean difference $\bar\Delta=\frac1n\sum_i(m_i^{\mathrm{ours}}-m_i^{\mathrm{ref}})$ and the one-sided 95\,\% lower bound $L_{95}$
 of its image-level bootstrap distribution, non-inferiority holds if
 \begin{equation}
 L_{95}(\bar\Delta)>-\delta,
 \label{eq:ni}
 \end{equation}
 with margin $\delta=0.5$~dB; SSIM has a margin of 0.01 and is reported either way; ours against the
 no-twin control uses a margin of 0.25~dB (the cost of the regulariser). Metrics are computed after a per-image gain and offset fit to the DES reference
 (``calibrated'', primary) and without it (``raw'').
\item[H3$'$, lesion absorption.] On VinDr-RibCXR sites, $A$(ours)$<A$(reference model) and $A$(ours)$<A$(no-twin control), as paired site-level
 differences with image-cluster bootstrap intervals, Holm over the two contrasts, and an interaction between sites that overlap an annotated rib and sites
 that do not.
\item[H4$'$, evidence without insertion.] (a) Contrast-to-noise change of real nodules on JSRT and Node21, ours against the reference model and against
 the no-twin control, for all lesions and for lesions with at least 15\,\% bone overlap. For a nodule box $\mathcal{B}$ and the ring
 $\mathcal{U}$ around it (the box grown by half its size on each side),
 \begin{equation}
\begin{gathered}
\mathrm{CNR}(I)=\frac{\mathrm{mean}_{\mathcal{B}}(I)-\mathrm{mean}_{\mathcal{U}}(I)}{\mathrm{sd}_{\mathcal{U}}(I)},\\
\Delta\mathrm{CNR}=\mathrm{CNR}(\hat s)-\mathrm{CNR}(x),
\end{gathered}
\label{eq:cnr}
\end{equation}
 compared between suppressors as paired differences per nodule. (b) Downstream: NIH ChestX-ray14 macro AUROC per input
 (original, ours, reference model, CLAHE, original plus ours; three seeds); TBX11K, Node21 and VinBigData detection by FROC and mAP@50; recall at two false
 positives per image for lesions with and without bone overlap, as a difference in differences against the original input.
\item[H5$'$, reliability.] Spearman correlation between ensemble spread and absolute bone error on the SZCH held-out pairs (pseudo-truth bone from the
 least-squares-aligned DES image); split-conformal coverage calibrated on Qure-projection validation CTs and measured on the Qure-projection test split
 (nominal 90\,\%) and, as measured under shift, on SZCH.
\end{description}

\subsection{Metrics and downstream models}
Paired fidelity uses PSNR and SSIM \citep{wang2004ssim} on the bone region and on the whole image, band-limited SSIM, MS-SSIM \citep{wang2003msssim}, LPIPS \citep{zhang2018lpips}
and a high-frequency energy ratio; the same smoothing is applied to every method and never tuned. The rib benchmark reports the relative reduction of
Sobel gradient energy on the rib-boundary band inside the lungs and on lung pixels away from ribs, their difference (selectivity) and the ROC-AUC of the
gradient-reduction map. Detection is scored by FROC sensitivity at fixed false-positive rates, averaged as the competition performance metric
\citep{vanginneken2010anode09,niemeijer2011cpm}, and by mAP@50; AUCs are compared with the paired DeLong test \citep{delong1988auc}. The downstream classifier is a DenseNet-121
\citep{huang2017densenet} and the detector a Faster R-CNN with a ResNet-50 FPN backbone \citep{ren2017fasterrcnn}, with identical settings for every
input variant (\path{s4/cls.py}, \path{s4/det.py}). The classifier sees $256$~px images, is trained for 10 epochs with AdamW
\citep{loshchilov2019adamw} (learning rate $3\times10^{-4}$, weight decay $10^{-4}$, batch 64, cosine schedule) and keeps the epoch with the
best validation macro-AUROC. The detector sees $512$~px images, is trained for 12 epochs with SGD (learning rate 0.01, momentum 0.9, weight
decay $10^{-4}$, batch 8) and keeps the epoch with the best validation CPM. A multi-channel input (original plus soft, or original, soft and bone)
stacks the images as channels. Detection metrics treat tied scores as one threshold, so they do not depend on the order in which a sort
returns ties.

\subsection{Encoder screen}
\label{sec:encoders}
To check whether a newer encoder helps under the same recipe, ConvNeXt V2 \citep{woo2023convnextv2}, DINOv3 \citep{simeoni2025dinov3} and RAD-DINO
\citep{perezgarcia2025raddino} were trained at a reduced budget under the same recipe, and an encoder moves to the full budget only if
it meets a rule fixed in advance (RUNPLAN.md, rule R1): at \num{100000} samples, absorption $A\le0.070$ on the development audit with
bone-region PSNR$_{cal}\ge21.60$~dB on the development pairs, or PSNR$_{cal}\ge21.95$~dB with $A\le0.100$. The thresholds come from the
SwinV2-T range at that budget and were written after one SwinV2-B run had been seen; the decision at the full budget used rule R2 (see Models and comparators above). Pretrained encoders other than SwinV2 are trained with a learning rate $0.25\times$ that of the decoder;
ConvNeXt-V2-T diverged at the full rate (fp16 overflow) and was rerun at $0.25\times$, an extension of the rule written for vision transformers.
DINOv3 ViT-S/16 did not improve on SwinV2-T, RAD-DINO ViT-B/14 missed R1 on PSNR by 0.05~dB but was continued at the full budget as an
exploratory encoder, and ConvNeXt-V2-T met R1.

\subsection{External validation protocol (BDCXR-3257)}
\label{sec:bdcxr_protocol}
An addendum to the analysis plan (SAP\_bdcxr.md) was frozen before any model of this study saw a BDCXR-3257 image; its one change, the
duplicate grouping (Section~\ref{sec:data}), was made from the images alone before any model ran and is recorded in the deviations log. All
suppressors are the frozen models; only the downstream classifiers are trained. (E1, primary) Pneumonia classification with three one-vs-rest
outputs, inputs as in H4$'$b (original, ours, reference model, CLAHE, original plus ours), the classifier settings of the downstream protocol,
five folds stratified by class over the near-duplicate groups (fold seed 0) with 10\,\% of each training part, by whole groups, for epoch selection,
and three seeds per fold and input (75 runs). Out-of-fold probabilities, averaged over seeds, give one prediction per image. The primary measure
is the pooled macro-AUROC; ours against the original image is non-inferior if the lower bound of the 95\,\% interval of the difference (group
bootstrap) exceeds $-0.01$, and superior if it exceeds 0; per class, paired DeLong tests with Holm over the three classes. (E2, secondary) The
lesion audit of Section~\ref{sec:audit} on a fixed class-stratified subset of 300 images (seed 7), six sites per image, with bone overlap from the
reference model's bone map; images whose lung mask covers less than 2\,\% of the image are skipped and counted (none was). (E3, descriptive) The
per-image ensemble spread over the lung field, compared with the held-out SZCH pairs and correlated with age.
